\honest{Privileging the Hypothesis}{Eliezer Yudkowsky}{2009}

A murder in a town of a million leaves no clues. A detective admits there is no evidence against anyone, and then proposes that we consider Mortimer Q. Snodgrass. I call this ``the fallacy of privileging the hypothesis'', and I ask whether it already has a name.

The detective needs some evidence in hand to single Mortimer out, though not evidence a court would accept. Without it, even asking the police to think about Mortimer violates Mortimer's rights. Suppose three detectives each name an enemy, and a witness then reports a brown-haired person. Mortimer is the brown-haired one, so the police decide Mortimer did it.

The reason is arithmetic. To narrow a billion boxes down to ten takes 27 bits of evidence; to pick the right one of the ten with better than even odds takes only 4 more. \nb{The figures check. Yudkowsky had made the same count in ``Einstein's Arrogance'' in 2007.} Naming one person out of a million with no reason skips most of the evidence.

Now quantum mechanics. In a recent conversation, Scott Aaronson admitted that no concrete evidence favors a collapse postulate or a single world, but pointed to possible future evidence and to the open question of the Born probabilities. This is privileging the hypothesis. ``There must be a trillion better ways to answer the Born question'' than a collapse that breaks every property of physical law. I give no example. \nb{The one answer Yudkowsky had discussed in 2008, Robin Hanson's, was put ``under fifty percent''.} \nb{Nothing singles out Mortimer. But the evidence that picked out quantum mechanics picked out the collapse and no-collapse readings together, since they predict the same results so far. So the charge depends on the premise that collapse is so complex that its rivals are ``no more complicated or unlikely'': the simplicity argument for many-worlds. In the published transcript, Aaronson disputed that premise, saying that one world adds less complexity than Yudkowsky's alternatives. The post does not report this.}

Whatever people think about, they treat as ``in the running''. I give no evidence. \nb{Research supports it: in studies reviewed by Koehler in 1991, people who imagined a possibility became more confident that it was true.} The theist who takes gaps in science as evidence for Jehovah, rather than for any of a trillion other gods, makes the same mistake, and so does the advocate of intelligent design. \nb{The first half of this is the old criticism of the ``God of the gaps''. The post adds the demand for evidence that points to one god in particular.} Complicated ideas should need more work to be raised to attention, as they need more evidence to be believed. In large spaces, ``attention without evidence is more than halfway to belief without evidence.''

A creationist moves doubt within evolutionary theory onto design. In the same way, doubt about the Born statistics should stay among theories that are continuous, linear, local and causal. But ``single-world theories violate all these characteristics''. \nb{Bell's theorem does make every single-world account nonlocal. But Bohmian mechanics, a single-world theory, keeps the linear, unitary evolution of the wave function and moves particles deterministically and continuously.}
